\section*{अध्याय \ref{ch:b3:learning-memory} --- तंत्रिकीय सुघट्यता, अधिगम और स्मृति}

\begin{solution}{exo:b3:learning-memory:1}
H.M. में \omterm{def:b3:learning-memory:kinds}{घोषणात्मक स्मृति} (तथ्य, घटनाएँ) मिट गई थी --- उसे कुछ भी नया
याद नहीं आता था --- जबकि अघोषणात्मक स्मृति बची रही: उसने आईने में
चित्र बनाना सामान्य ढंग से सीखा, बिना उन सत्रों को याद रखे। कौशल:
धारीदार पिंड और \omterm{def:b3:nervous-systems:plan}{अनुमस्तिष्क}; अनुबंधित भय: अमिग्डला; अनुबंधित प्रेरक
प्रतिवर्त: \omterm{def:b3:nervous-systems:plan}{अनुमस्तिष्क}; पूर्वप्रभावन: संवेदी प्रांतस्था।
\end{solution}

\begin{solution}{exo:b3:learning-memory:2}
जब कोई पूर्वसंधिक कोशिका बार-बार किसी पश्चसंधिक कोशिका को दगाने में
योग दे, तब उनका जोड़ मज़बूत होता है। इसीलिए LTP निवेश-विशिष्ट है (केवल
वे सक्रिय संधियाँ बदलती हैं जिन्होंने हिस्सा लिया), साहचर्यी है (कोई
दुर्बल निवेश, जो तब सक्रिय हो जब वह कोशिका किसी प्रबल के तहत दग रही
हो, मज़बूत होता है), और सहकारी है (उस कोशिका को दगाने के लिए पर्याप्त
निवेशों को साथ काम करना पड़ता है)।
\end{solution}

\begin{solution}{exo:b3:learning-memory:3}
वह केवल तब चालन करता है जब ग्लूटामेट बँधा हो (यानी पूर्वसंधिक कोशिका
दगी) \emph{और} झिल्ली इतनी विध्रुवित हो कि मैग्नीशियम निकल जाए (यानी
पश्चसंधिक कोशिका सक्रिय है): यानी हेब के नियम की दोनों शर्तें किसी एक
ही प्रोटीन में। टेटेनस से पहले AP5: कोई कैल्सियम-प्रवेश नहीं, कोई LTP
नहीं, हालाँकि \omterm{prop:b3:learning-memory:nmda}{AMPA ग्राहियों} से होता संचरण सामान्य है। बाद में AP5:
प्रेरित हो चुका LTP अछूता रहता है --- वह ग्राही प्रेरण के लिए चाहिए,
रख-रखाव के लिए नहीं।
\end{solution}

\begin{solution}{exo:b3:learning-memory:4}
अल्पकालिक: पहले से मौजूद प्रोटीनों का सहसंयोजी रूपांतरण --- PKA किसी
पोटैशियम वाहिका का फ़ॉस्फ़ोरिलन करके उसे बंद कर देती है, जिससे आवेग
चौड़ा होता है और मोचन बढ़ता है; कोई नया प्रोटीन नहीं। दीर्घकालिक: PKA
CREB सक्रिय करती है, नए जीन अनुलेखित होते हैं और नए संधि-सिरे बढ़ते
हैं। बार-बार के प्रशिक्षण के दौरान दिया गया कोई प्रोटीन-संश्लेषण निरोधक
अल्पकालिक सुकरण अछूता छोड़ देता है और दिनों बाद नापी गई स्मृति मिटा
देता है; बाद में दिया जाए तो वह कुछ नहीं करता।
\end{solution}

\begin{solution}{exo:b3:learning-memory:5}
$0.7^{n} \le 0.1$: $n \ge 6.5$, इसलिए $7$ आज़माइशें। $0.9\lambda$ से,
$0.9\times 0.7^{n} < 0.1$: फिर $7$ आज़माइशें। सममित इसलिए कि दोनों
दिशाओं को वही $\alpha$ चलाता है। प्राणियों में विलोपन प्रायः धीमा होता
है और मूल स्मृति विश्राम के बाद लौट आती है (स्वतःस्फूर्त
पुनःप्राप्ति): विलोपन कोई नया अधिगम है जो पुराने को दबाता है, उसका मिटाव
नहीं।
\end{solution}

\begin{solution}{exo:b3:learning-memory:6}
अभिलक्षणिक सदिश $(1,1)/\sqrt{2}$ अभिलक्षणिक मान $1.5$ के साथ और
$(1,-1)/\sqrt{2}$ $0.5$ के साथ; हेब $\mathbf{w}$ को $(1,1)$ की ओर
मोड़ता है, और वह न्यूरॉन दोनों निवेशों के साथ दगने को पकड़ने लगता है।
ओया का पग: $y = 1$,
$\Delta\mathbf{w} = 0.1\times 1\times((1,1) - 1\times(1,0)) = (0,0.1)$,
इसलिए $\mathbf{w} = (1, 0.1)$।
\end{solution}

\begin{solution}{exo:b3:learning-memory:7}
$500\times 0.05 = 25$ मुक्त आयन; $25/6\times 10^{23} = 4.2\times 10^{-23}$
मोल, $10^{-16}$ L में: \qty{0.42}{\micro mol/L},
यानी विश्रामी स्तर से चार गुना पर कैल्मॉड्युलिन के $K_{d}$ से नीचे ---
यानी केवल आंशिक सक्रियण; उस माइक्रोमोलर परास तक पहुँचने के लिए, जो
CaMKII चालू कर देती है, कुछ वाहिकाओं का साथ खुलना चाहिए।
\end{solution}

\begin{solution}{exo:b3:learning-memory:8}
A को तब तक प्रशिक्षित किया जाता है जब तक $V_{A} = \lambda$ न हो जाए:
डोपामीन आरंभिक आज़माइशों पर झोंके देता है और पुरस्कार के पूर्वानुमानित
हो जाने पर चुप हो जाता है। संयुक्त आज़माइशों पर पूर्वानुमान
$V_{A} + V_{B} = \lambda$ पहले से ही ठीक है, त्रुटि शून्य है, डोपामीन
चुप है, और $V_{B}$ कभी बदलता ही नहीं: B अवरुद्ध है। अकेले जाँचने पर B
कोई अनुक्रिया नहीं खींचता; और यदि तब कोई पुरस्कार अकेले B के बाद आए तो
वह अप्रत्याशित है और डोपामीन झोंका देता है।
\end{solution}

\begin{solution}{exo:b3:learning-memory:9}
(क) \omterm{def:b3:learning-memory:hippocampus}{CA1} में कोई LTP नहीं और कोई स्थानिक-स्मृति ह्रास, जबकि दूसरा अधिगम
अक्षत। (ख) अल्पकालिक स्मृति सामान्य, \qty{24}{h} पर दीर्घकालिक स्मृति
अनुपस्थित। (ग) कोई असर नहीं --- \omterm{def:b3:learning-memory:kinds}{सुदृढ़ीकरण} पूरा हो चुका है (जब तक वह
स्मृति फिर सक्रिय न की जाए, जिस पर वह दुबारा अस्थिर हो जाती है)। (घ) वह
चूहा सुरक्षित संदर्भ में जड़ हो जाता है: उन्हीं कोशिकाओं को फिर सक्रिय
करके, जिन्होंने उसे कूटबद्ध किया था, वह स्मृति कृत्रिम ढंग से याद करा
दी जाती है।
\end{solution}

\begin{solution}{exo:b3:learning-memory:10}
$\Delta|\mathbf{w}|^{2} \approx 2\mathbf{w}\cdot\Delta\mathbf{w} = 2\eta
y^{2} \ge 0$: वह लंबाई कभी घटती नहीं और जब भी वह न्यूरॉन दगे तब बढ़ती
है। ओया के किसी स्थिर बिंदु पर $C\mathbf{w} = (\mathbf{w}^{\mathsf{T}}C
\mathbf{w})\mathbf{w}$, इसलिए $\mathbf{w}$ कोई अभिलक्षणिक सदिश है जिसमें
$\lambda =
\lambda|\mathbf{w}|^{2}$, और इसलिए $|\mathbf{w}| = 1$। असीम वृद्धि हर
संधि संतृप्त कर देती, वरणशीलता मिटा देती और भगोड़ा उत्तेजन चला देती;
जैविक अभिसामान्यक संधि-मापक्रमण (किसी न्यूरॉन के बहुत सक्रिय होने पर
उसकी सारी संधियाँ नीचे मापक्रमित), LTD, और नींद में संधियों का
नीचे-मापक्रमण हैं।
\end{solution}

\begin{solution}{exo:b3:learning-memory:11}
चूहा: $0.14\times 3\times 10^{5} \approx 4\times 10^{4}$; मनुष्य: $0.14
\times 2\times 10^{6} \approx 3\times 10^{5}$। यह आकलन यादृच्छिक सघन
प्रतिरूप और पूरी संयोजकता मानता है; असली प्रतिरूप विरल (जो क्षमता बहुत
बढ़ा देता है) और संरचित होते हैं, और \omterm{def:b3:learning-memory:kinds}{हिप्पोकैम्पस} कोई अस्थायी भंडार है
जो स्मृतियाँ $10^{10}$ न्यूरॉनों की किसी प्रांतस्था को सौंप देता है ---
जीवन भर का भंडार वही प्रांतस्था है, और वह हिप्पोकैम्पसी संख्या कोई
बफ़र-आकार है, स्मृतियों की गिनती नहीं।
\end{solution}

\begin{solution}{exo:b3:learning-memory:12}
खुली आँख के निवेश उस प्रांतस्थीय कोशिका के दगने से सहसंबद्ध हैं और
मज़बूत होते हैं; बंद आँख के निवेश नहीं, और किसी सीमित कुल भार के लिए
स्पर्धा में वे कमज़ोर पड़ जाते हैं --- यानी हेब जमा अभिसामान्यन। दोनों
आँखें बंद हों तो किसी भी निवेश का पक्ष नहीं लिया जाता, इसलिए दोनों
अपना (दुर्बल) इलाक़ा रखे रहते हैं और असर हलका होता है। वह काल तब बंद
होता है जब निरोधी परिपथ परिपक्व होते हैं, संधियों के चारों ओर
परिन्यूरॉनी जाल कड़े पड़ते हैं, और \omterm{prop:b3:learning-memory:nmda}{NMDA ग्राही} की उपइकाइयाँ बदलकर किसी
कम उदार रूप में आ जाती हैं।
\end{solution}

\begin{solution}{pb:b3:learning-memory:1}
\textbf{1.} $10\times 10^{3} = 10^{4}$ आयन; \qty{2}{\%} मुक्त, यानी
$200$; $200/6\times 10^{23} = 3.3\times 10^{-22}$ मोल, $8\times 10^{-17}$ L
में: \qty{4.2}{\micro mol/L}।
\textbf{2.} हाँ, \qty{2}{\micro mol/L} की देहली से ऊपर। दो
ग्राही: $40$ मुक्त आयन, \qty{0.83}{\micro mol/L} --- यानी उससे नीचे।
\textbf{3.} \qty{0.83}{\micro mol/L} पर कैल्सिन्यूरिन सक्रिय है और
CaMKII नहीं: \omterm{prop:b3:learning-memory:nmda}{AMPA ग्राही} हटा दिए जाते हैं और वह संधि कमज़ोर होती है
(LTD)। एक आयन, भिन्न बंधुताओं की दो एंज़ाइम: कोई बड़ी तेज़ बढ़त उस
काइनेज़ तक पहुँचती है, कोई छोटी धीमी केवल उस फ़ॉस्फ़ेटेज़ तक।
\textbf{4.} कुछ ही काल-अचरों के भीतर आया कोई दूसरा निवेश अपना कैल्सियम
पहले के बचे हुए में जोड़ देता है; \qty{50}{ms} के बाद लगभग कुछ भी नहीं
बचता। NMDA की संपात-शर्त और \qty{20}{ms} का क्षय मिलकर दोनों ओर कोई
\qty{20}{ms} की खिड़की देते हैं --- यानी वही आवेग-काल की खिड़की।
\textbf{5.} बफ़र अधिकांश कैल्सियम माइक्रोमीटर के किसी अंश के भीतर पकड़
लेते हैं और उस कंटक की सँकरी गर्दन द्रुमिका तक विसरण धीमा कर देती है,
इसलिए वह बढ़त उसी कंटक में रहती है जो सक्रिय था; \qty{0.5}{\micro m}
दूर पड़ा कोई पड़ोसी कुछ नहीं देखता, और केवल वही संधि बदलती है जो सक्रिय
थी।
\textbf{6.} $(70 - 30)/5 = 8$ समकालिक निवेश, या कोशिका-काया से पीछे
लौटता कोई क्रिया विभव: कोई अकेली संधि अपने ही \omterm{prop:b3:learning-memory:nmda}{NMDA ग्राही} नहीं खुलवा
सकती --- \omterm{def:b3:structural-biology:allostery}{सहकारिता} भौतिकी में ही बनी हुई है।
\textbf{7.} $(1,1)/\sqrt{2}$, अभिलक्षणिक मान $1.6$ के साथ;
$(1,-1)/\sqrt{2}$, $0.4$ के साथ।
\textbf{8.} $C\mathbf{w}_{0} = (1, 0.6)$, $\mathbf{w}_{1} = (1.10,
0.06)$, $(1,1)$ से कोण: $45^{\circ} - 3.1^{\circ} = 41.9^{\circ}$।
$\mathbf{w}_{2} = (1.21, 0.13)$: $38.8^{\circ}$। $\mathbf{w}_{3} = (1.34,
0.22)$: $35.8^{\circ}$।
\textbf{9.} $w_{2}/w_{1} \to 1$, जबकि $|\mathbf{w}|$ बिना किसी सीमा के
बढ़ता है (प्रति पग गुणक $1.16$ से)।
\textbf{10.} $C\mathbf{w} = (1.28, 1.28)$; $\mathbf{w}^{\mathsf{T}}C
\mathbf{w} = 2.05$; हर एक के लिए
$\Delta\mathbf{w} = 0.1\times(1.28 - 2.05\times 0.8)
= -0.036$: $\mathbf{w} = (0.76, 0.76)$, जो इकाई सदिश
$(0.71, 0.71)$ की ओर सिकुड़ रहा है।
\textbf{11.} संयुक्त $y = 1.42$; अकेला निवेश 2 $y = 0.71$, यानी आधा।
\textbf{12.} वह दुर्बल निवेश तब दगता है जब प्रबल उस कोशिका को चला रहा
हो, इसलिए वह निर्गम से सहसंबद्ध है और \omterm{def:b3:learning-memory:ltp}{हेब का नियम} उसे मज़बूत कर देता है
--- वह न्यूरॉन उस सह-घटना को सीख लेता है।
\textbf{13.} $1 - 0.85^{n}$: $0.56$, $0.80$, $0.96$।
\textbf{14.} $0.85^{n} \le 0.05$: $n \ge 18.5$, इसलिए $19$ आज़माइशें।
\textbf{15.} नहीं: $n$ आज़माइशों के बाद छुआ गया अनंतस्पर्शी का अंश
$\lambda$ पर निर्भर नहीं करता; वह अनंतस्पर्शी दुगुना हो जाता है। कोई
बड़ा पुरस्कार $\alpha$ बढ़ा सकता है (प्रमुखता), और
\emph{निरपेक्ष} मूल्य पर टिकी कोई व्यवहारिक देहली जल्दी छू ली जाती है।
\textbf{16.} $V_{B} = 0$: अवरुद्ध। वह संयुक्त पूरी तरह पूर्वानुमानित है,
त्रुटि शून्य है, डोपामीन न्यूरॉन चुप हैं।
\textbf{17.} अकेला B कुछ भी पूर्वानुमानित नहीं करता, इसलिए त्रुटि
$\lambda$ है और डोपामीन झोंका देता है; फिर $V_{B}$ अर्जन की तरह चढ़ता
है, $\lambda(1 -
0.85^{n})$।
\textbf{18.} वह औषध वही झोंका पहुँचा देती है जिसका अर्थ है
``पूर्वानुमान से बेहतर'', चाहे उससे पहले कुछ भी हुआ हो; वह नियम हर मौके
पर उन्हीं संकेतों तथा क्रियाओं का मूल्य बढ़ा देता है, इसलिए वे किसी भी
असली परिणाम से स्वतंत्र होकर मूल्य पा लेते हैं --- और वह खोज बाध्यकारी
बन जाती है।
\textbf{19.} $0.14\times 3\times 10^{5} \approx 4.2\times 10^{4}$।
\textbf{20.} $4.2\times 10^{4}/20 = 2100$ दिन, यानी लगभग छह वर्ष;
\omterm{def:b3:learning-memory:kinds}{सुदृढ़ीकरण} को स्मृतियाँ प्रांतस्था तक ले जानी होंगी और विस्मरण को वह
भंडार खाली करना होगा।
\textbf{21.} नए प्रतिरूप जल्दी-जल्दी प्रांतस्था में लिखने से वे संधियाँ
मिट जातीं जो पुरानों को कूटबद्ध करती हैं (विनाशकारी दख़ल); धीमा,
गुँथा हुआ \omterm{def:b3:learning-memory:hippocampus}{पुनर्चालन} प्रांतस्था को नए को पुराने के साथ समाकलित करने देता
है, जबकि तब तक वह नई स्मृति \omterm{def:b3:learning-memory:kinds}{हिप्पोकैम्पस} थामे रहता है --- यानी कोई तेज़
बफ़र, जो किसी धीमे भंडार को खिलाता है।
\textbf{22.} $0.03\times 3\times 10^{5} = 9000$ कोशिकाएँ प्रति स्थान।
विरल प्रतिरूप कम अतिव्यापी होते हैं, इसलिए उनका दख़ल पुनःप्राप्ति
बिगाड़े उससे पहले कहीं अधिक संचित किए जा सकते हैं; प्रतिरूपों के विरल
होते जाने पर क्षमता $N$ से भी तेज़ी से चढ़ती है, और कीमत यह कि हर
प्रतिरूप विशिष्टता के लिए अधिक कोशिकाएँ लगाता है।
\textbf{23.} $20$ का निचोड़: जो कोशिकाएँ उस दौड़ में \qty{400}{ms} की
दूरी पर दगीं वे \omterm{def:b3:learning-memory:hippocampus}{पुनर्चालन} में \qty{20}{ms} की दूरी पर दगती हैं --- यानी
आवेग-काल की खिड़की के भीतर। उस निचोड़ का मतलब यही है कि जो क्रम असली
समय में जुड़ने के लिए बहुत धीमे हैं वे हेबीय पहुँच में आ जाएँ।
\textbf{24.} एक दिन पर वह स्मृति अब भी \omterm{def:b3:learning-memory:kinds}{हिप्पोकैम्पस} पर टिकी है; महीने
भर पर वह प्रांतस्था में सुदृढ़ हो चुकी है और उसे उसकी ज़रूरत नहीं रही
--- यानी ठीक H.M. का काल-श्रेणीबद्ध लोप, जिसमें हाल की स्मृतियाँ गईं और
दूर की बची रहीं।
\textbf{25.} दस खुलनों के बाद मुक्त कैल्सियम लगभग
\qty{4.2}{\micro mol/L}; \qty{95}{\%} तक $19$ आज़माइशें; और चूहे के
\omterm{def:b3:learning-memory:hippocampus}{CA3} में कोई $4\times 10^{4}$ प्रतिरूप।
\end{solution}
